\begin{table}[htbp]
\centering
\begin{threeparttable}
\caption{Residual-asymmetry and return-only benchmark comparisons.}
\label{tab:secondary-comparisons}
\scriptsize
\setlength{\tabcolsep}{4pt}
\begin{tabular}{@{}llrrr@{}}
\toprule
Comparison & Horizon & S\&P~500 & FTSE~100 & DAX \\
\midrule
\multicolumn{5}{l}{\textit{Panel A: RV-RA-NN-SV $-$ RV-NN-SV}} \\
@@ROW_0@@
@@ROW_1@@
@@ROW_2@@
\addlinespace
\multicolumn{5}{l}{\textit{Panel B: NN-SV $-$ GARCH$(1,1)$-$t$}} \\
@@ROW_3@@
@@ROW_4@@
@@ROW_5@@
@@ROW_6@@
\bottomrule
\end{tabular}
\begin{tablenotes}[flushleft]
\footnotesize
\item Notes: Entries are mean loss differences, with two-sided descriptive DM $p$-values in brackets. Negative entries favour the first model named in each panel. RV-RA comparisons use pairwise available samples; the smallest common sample contains 899 origins. The GARCH$(1,1)$-$t$ comparisons use the corresponding common samples reported for the main matched comparisons in Table~\ref{main-tab:matched-comparisons}. GARCH-$t$ is supplementary and does not enter the nine-model confidence sets.
\end{tablenotes}
\end{threeparttable}
\end{table}
